\section{HeatGeo}
\label{sec:method}

\subsection{Teacher Semantic Graph}

HeatGeo first encodes the full training corpus with the teacher and constructs a teacher semantic graph. For each example $x_i$, let
\begin{equation}
    \mathcal{N}_k^T(i)
    =
    \operatorname{TopK}_{j\neq i}\cos(t_i,t_j)
\end{equation}
be its top-$k$ teacher nearest neighbors. To reduce noisy one-directional edges and hubness, HeatGeo uses a mutual $k$NN graph:
\begin{equation}
    j \in \mathcal{M}_k^T(i)
    \Longleftrightarrow
    j \in \mathcal{N}_k^T(i)
    \;\mathrm{and}\;
    i \in \mathcal{N}_k^T(j).
\end{equation}
The weighted adjacency matrix $W^T$ is
\begin{equation}
    W_{ij}^{T}
    =
    \begin{cases}
    \exp(\cos(t_i,t_j)/\tau_g), & j \in \mathcal{M}_k^T(i),\\
    0, & \mathrm{otherwise},
    \end{cases}
\end{equation}
where $\tau_g$ controls graph sharpness. Let $D^T$ be the diagonal degree matrix. The row-normalized transition matrix is
\begin{equation}
    P^T = (D^T)^{-1}W^T.
\end{equation}
Each row $P_i^T$ defines a one-step semantic random-walk distribution from anchor $x_i$.

The graph temperature $\tau_g$ is not a cosmetic hyperparameter. A mutual $k$NN graph has degree of order $k$, and teacher cosine similarities among the top-$k$ neighbors of an anchor occupy a narrow band. If $\tau_g$ is large relative to that band, every row of $P^T$ is close to uniform over its support, and the diffusion targets below degrade into binary neighbor/non-neighbor labels that carry no ranking information. We therefore report $D_{\mathrm{KL}}(p_{i,r}^T\|\mathcal{U}_{\mathrm{supp}})$, the divergence of each target from the uniform distribution on its own support, as a build-time diagnostic and select $\tau_g$ so that this quantity stays well above zero.

\subsection{Multi-Scale Heat Diffusion Targets}

A one-step transition captures direct neighbors, but embedding spaces also contain higher-order semantic structure. HeatGeo therefore uses multiple diffusion scales of a \emph{lazy} random walk,
\begin{equation}
    Q^T = \tfrac{1}{2}\left(I + P^T\right),
    \qquad
    P_r^T = (Q^T)^r,
    \qquad r \in \mathcal{R},
\end{equation}
where $\mathcal{R}$ is typically $\{1,2,4\}$. Small $r$ preserves sharp local neighborhoods, while larger $r$ captures semantic regions and multi-hop communities. Laziness is what makes $\mathcal{R}$ a graded family: a plain walk on a mutual $k$NN graph mixes within its local component in about two steps, so $(P^T)^2$ and $(P^T)^4$ collapse onto nearly the same distribution and the additional scales contribute duplicated terms rather than new structure. The lazy walk retains mass at intermediate step counts and yields scales that remain distinguishable. In implementation, diffusion is approximated sparsely by propagating probability mass through graph neighbor lists and retaining the top candidates, with a single walk snapshotted at each scale in $\mathcal{R}$. The truncation to a sparse support and the removal of the self-transition mass both destroy normalization, so the distribution is renormalized after each; we log the residual mass falling outside the retained support so that the approximation error is visible rather than assumed small.

For each anchor $x_i$, HeatGeo precomputes a diffusion \emph{pool} $\mathcal{P}_i$ of size $P$ --- the retained sparse support of the walk --- together with a same-source hard-negative pool $\mathcal{H}_i$ drawn from the teacher's top-ranked non-neighbors of $x_i$. At the start of each epoch $e$ it then \emph{redraws} a candidate set
\begin{equation}
    C_i^{(e)} = C_i^{+} \cup C_i^{h} \cup C_i^{r},
    \qquad |C_i^{(e)}| = C \ll P,
\end{equation}
where $C_i^{+}\subset\mathcal{P}_i$ contains diffusion neighbors, $C_i^{h}\subset\mathcal{H}_i$ contains teacher hard negatives, and $C_i^{r}$ contains random negatives, each filling a fixed slot quota. Resampling matters more than it appears. A candidate set frozen for the whole run defines a fixed finite collection of comparisons; once the student orders them correctly there is no further gradient, and training saturates after roughly one epoch regardless of how much structure the graph still holds. Redrawing $C_i^{(e)}$ exposes a different section of the same pool each epoch, so the objective keeps presenting comparisons the student has not yet fit while the underlying target distribution is unchanged. Within each scale we keep its top-$m$ neighbors deterministically and draw the remainder of its share by Gumbel top-$k$ sampling proportional to diffusion mass, which reaches the tail of the support without biasing the retained head.

Diffusion positives are drawn round-robin across the scales in $\mathcal{R}$, each scale sampling from its own $p_{i,r}^T$. Drawing the whole positive quota from the mixture $\sum_r\omega_r p_{i,r}^T$ is the more obvious choice and is wrong: the mixture's support is dominated by the broadest walk, so the sharpest scale receives one or two scored candidates and its KL term stops carrying ranking information at all. The teacher diffusion target at scale $r$ is the diffusion mass restricted and renormalized over $C_i$:
\begin{equation}
    p_{i,r}^{T}(j)
    =
    \frac{(P_r^T)_{ij}}
    {\sum_{j' \in C_i}(P_r^T)_{ij'}},
    \qquad j\in C_i.
\end{equation}
This target is a soft distribution over semantic relatedness, rather than a binary positive-negative label.

\subsection{Student Neighborhood Distribution}

Given an anchor embedding $s_i$ and candidate embeddings $\{s_j:j\in C_i\}$, the student induces a distribution over the same candidate set. A single student distribution, however, cannot support a multi-scale objective. Suppose the student used one temperature $\tau_s$ for all scales, giving one $p_i^S$. Since cross-entropy is linear in its target,
\begin{equation}
    \sum_{r\in\mathcal{R}}\omega_r D_{\mathrm{KL}}\!\left(p_{i,r}^{T}\,\|\,p_i^{S}\right)
    =
    \underbrace{-\sum_{r}\omega_r H\!\left(p_{i,r}^{T}\right)}_{\text{constant}}
    +\,
    H\!\left(\bar p_i^{T},\,p_i^{S}\right),
    \label{eq:collapse}
\end{equation}
where $H(\cdot,\cdot)$ is cross-entropy and $\bar p_i^{T}=\sum_{r}\omega_r p_{i,r}^{T}$. Up to an additive constant the objective depends on $\mathcal{R}$ only through the single mixture $\bar p_i^T$: any two scale sets with the same mixture produce identical gradients, and ``multi-scale'' reduces to target smoothing. Worse, \eqref{eq:collapse} is minimized at $p_i^S=\bar p_i^T$, at which point the loss equals
\begin{equation}
    \mathrm{JS}_{\omega}\!\left(p_{i,1}^{T},\dots,p_{i,R}^{T}\right)
    =
    H\!\left(\bar p_i^{T}\right)-\sum_{r}\omega_r H\!\left(p_{i,r}^{T}\right)
    \;\ge\;0,
    \label{eq:floor}
\end{equation}
the generalized Jensen--Shannon divergence of the scales. This is an \emph{irreducible floor}: it is a property of the teacher graph, fixed before training begins, and no student can go below it. A tied-temperature run that plateaus is therefore ambiguous --- the plateau may be an optimization failure or may simply be \eqref{eq:floor} --- and the two have opposite fixes. We compute \eqref{eq:floor} offline from the artifact and report the \emph{excess} $\mathcal{L}_{\mathrm{diff}}-\mathrm{JS}_\omega$ alongside the raw loss, so the distinction is observable at every step.

HeatGeo therefore assigns each scale its own student temperature $\tau_r$, with broader scales receiving larger $\tau_r$ to match their smoother targets:
\begin{equation}
    p_{i,r}^{S}(j)
    =
    \frac{\exp(\cos(s_i,s_j)/\tau_r)}
    {\sum_{j'\in C_i}\exp(\cos(s_i,s_{j'})/\tau_r)},
    \qquad j\in C_i,
\end{equation}
\begin{equation}
    \mathcal{L}_{\mathrm{diff}}
    =
    \frac{1}{|B|}
    \sum_{i\in B}
    \sum_{r\in\bar{\mathcal{R}}}
    \omega_r
    D_{\mathrm{KL}}
    \left(
        p_{i,r}^{T} \,\|\, p_{i,r}^{S}
    \right),
\end{equation}
where $B$ is a mini-batch, $\omega_r$ is the scale weight, and $\bar{\mathcal{R}}=\{0\}\cup\mathcal{R}$ includes the direct scale introduced in Section~\ref{sec:direct}. The scales now share the underlying similarity vector $\cos(s_i,\cdot)$ but not the distribution built from it, so \eqref{eq:collapse} no longer applies and \eqref{eq:floor} becomes a lower bound rather than an attained value. The mechanism is over-determination: $|\mathcal{R}|$ softmaxes at distinct temperatures impose $|\mathcal{R}|$ independent constraints on one $|C_i|$-vector of cosines, and satisfying them all requires reproducing the teacher's similarity \emph{gaps}, not merely its ranking. A single softmax pins $\cos(s_i,s_j)$ only up to a per-anchor additive constant, which is exactly the calibration that Spearman correlation on STS and any globally thresholded cosine retrieval depend on. In a controlled setting with free logits, minimizing the tied objective converges to the floor of \eqref{eq:floor} to four decimals, whereas untying $\tau_r$ reaches a value well below it.

\subsection{The Direct Scale}
\label{sec:direct}

The diffusion targets are graph mass renormalized over the scored columns, so every column outside $\mathcal{P}_i$ receives target \emph{exactly zero}. A zero target is not an abstention. The gradient of the cross-entropy with respect to the logit of column $j$ is $p_{i,r}^S(j)-p_{i,r}^T(j)$, so at a zero-target column it is $+p^S_{i,r}(j)>0$, and the objective drives $\cos(s_i,s_j)$ down without bound. Under in-batch sharing roughly $97\%$ of the columns an anchor is scored against are zero-target, and they include the hard negatives $C_i^h$ --- same-source, inside the teacher's top-200 by cosine, excluded from the graph only because the mutual-$k$NN test rejected them. In a controlled probe with planted cluster structure, those hard negatives carry teacher cosine $0.283$ against $0.349$ for the diffusion positives, and the graph-only objective pushes $98\%$ of them apart. The method is thus training the student to separate pairs its own teacher calls similar, which is precisely the calibration that STS Spearman and a global cosine threshold measure.

We therefore add a scale $r=0$ whose target is the teacher's own similarity over the full column set,
\begin{equation}
    p_{i,0}^{T}(j)
    =
    \frac{\exp\!\left(\cos(t_i,t_j)/\tau_t\right)}
    {\sum_{j'}\exp\!\left(\cos(t_i,t_{j'})/\tau_t\right)},
    \label{eq:direct}
\end{equation}
and optimize over $\bar{\mathcal{R}}=\{0\}\cup\mathcal{R}$. Unlike the diffusion targets, \eqref{eq:direct} is dense: every scored column receives its true teacher mass. In the same probe it assigns $0.670$ to the positives, $0.271$ to the hard negatives and $0.059$ to the random ones, and the fraction of hard negatives pushed apart falls from $98\%$ to $41\%$. The teacher embeddings are cached already, so the term is free.

Scale $r=0$ also supersedes $\mathcal{L}_{\mathrm{anchor}}$ on its own terms. Comparing cosines measured in the student's metric against cosines measured in the teacher's is \emph{not} invariant to a linear map of the student space --- an anisotropic rescaling of $s$ changes $\cos(s_i,s_j)$ and therefore changes \eqref{eq:direct}'s loss --- so this is the term that actually pins the absolute geometry, which Section~\ref{sec:objective} shows the projected anchor cannot. Note also that $\tau_t$ cannot simply inherit $\tau_g$: the graph temperature acts on a mutual-$k$NN neighborhood where cosines occupy a narrow band, whereas \eqref{eq:direct} ranges over the whole column set, where the same value would place nearly all mass on the single nearest column.

Finally, we score every anchor in $B$ against the \emph{union} of all candidates in the batch rather than against its own $C_i$ alone. The candidate texts are already encoded, so this is free in forward cost, multiplies the negatives per anchor by roughly $|B|$, and --- because anchors now share columns --- couples their logit scales so that similarity levels become comparable across the batch. The teacher target is scattered onto the union with zero mass on the columns an anchor did not draw, and the anchor's own column is masked.

\subsection{Why No Separate Spectral Term}
\label{sec:no-spectral}

A natural addition would be a spectral term that regresses student embeddings onto the leading eigenvectors of the graph Laplacian $L^T = I - (D^T)^{-1/2}W^T(D^T)^{-1/2}$, on the grounds that diffusion matching is local while eigenvectors describe global topology. The usual argument for omitting it appeals to the spectral expansion of the heat kernel,
\begin{equation}
    e^{-tL^T} = \sum_{\ell} e^{-t\lambda_\ell}\, u_\ell u_\ell^{\top},
    \label{eq:heatspec}
\end{equation}
concluding that matching diffusion at several scales already matches the graph spectrum under a low-pass filter of bandwidth $t$, so an eigenvector term distills the same structure twice.

That argument does not survive Section~\ref{sec:no-spectral}'s own premises if the temperatures are tied. By \eqref{eq:collapse}, a tied objective constrains the student only through the mixture $\bar p^T$, i.e.\ through the \emph{single} fixed filter $g(Q^T)=\sum_r\omega_r (Q^T)^r$ applied once. Observing $u_\ell^\top g(Q^T) u_\ell = g(1-\lambda_\ell)$ for one known $g$ does not identify $\{\lambda_\ell\}$: $g$ is not injective on $[-1,1]$ for the weights we use, and in any case one scalar functional per eigenvalue recovered from row-normalized, candidate-restricted distributions is far weaker than eigenvector regression. Under a tied temperature the ``already covered'' claim is therefore unsupported, and the honest statement would be that the two terms are complementary.

With per-scale temperatures the argument is restored, in a weaker and more defensible form. Each scale now imposes its own constraint, so the student must match $|\mathcal{R}|$ distinct filters $(Q^T)^r$ of the same operator rather than one mixture --- a filter bank, and it is the bank, not any single member, that pins down where the spectral mass sits. We still make no claim of exact spectral identifiability from $|\mathcal{R}|=3$ scales on truncated supports; we claim only that a truncated eigenvector-regression term adds little on top of the bank while reintroducing the pipeline's only global eigendecomposition. Section~\ref{sec:experiments} reports the ablation that adds the term back, run against the per-scale-temperature model so that the comparison is against the objective we actually advocate.

\subsection{Training Objective}
\label{sec:objective}

A relational objective constrains the student only up to a transformation of its space, which invites a pointwise anchor of the form
\begin{equation}
    \mathcal{L}_{\mathrm{anchor}}
    =
    \frac{1}{|B|}
    \sum_{i\in B}
    \left(1-\cos(W_a s_i,t_i)\right),
    \label{eq:anchor}
\end{equation}
where $W_a$ is a freely learned map to the teacher dimension. We report this term but disable it by default ($\lambda_{\mathrm{anchor}}=0$), because as written it cannot do the job it is introduced for. For any invertible $A$, replacing $s_i\mapsto As_i$ and $W_a\mapsto W_aA^{-1}$ leaves \eqref{eq:anchor} exactly invariant, so the term is blind to precisely the degrees of freedom --- rotation, anisotropic scaling, shear --- that the relational loss also leaves free. Its optimum is the best linear predictor of $t_i$ from $s_i$, a quantity that a single linear layer can fit early and then hold nearly constant; empirically it settles within one epoch and contributes no further gradient of consequence. Supplying genuine absolute calibration requires either freezing $W_a$ (e.g.\ to a random projection, breaking the gauge) or, as we do here, obtaining it from the per-scale temperatures, which constrain cosine gaps directly in the student's own metric. We keep \eqref{eq:anchor} only as an ablation row.

The final objective is therefore two terms with a single balancing weight:
\begin{equation}
    \mathcal{L}_{\mathrm{HeatGeo}}
    =
    \mathcal{L}_{\mathrm{diff}}
    +
    \lambda_{\mathrm{anchor}}\mathcal{L}_{\mathrm{anchor}}.
\end{equation}
All scales in $\mathcal{R}$ are active from the first step. We also evaluated a curriculum that introduces scales one epoch at a time; because the scale weights are renormalized over the active set, such a curriculum monotonically shifts weight away from the sharpest local scale as training proceeds, and we found no benefit from it (Section~\ref{sec:experiments}).

The loss value alone is not interpretable. It cannot separate a student that has learned the teacher's neighborhood ranking from one whose embeddings have collapsed toward a uniform candidate distribution, and by \eqref{eq:floor} it cannot separate a converged student from an exhausted objective. We therefore log, at every step, (i) the entropy of $p_{i,r}^S$ relative to $\log|C_i|$, (ii) the student probability mass falling on the teacher's top-8 candidates, and (iii) the excess $\mathcal{L}_{\mathrm{diff}}-\mathrm{JS}_\omega$ over the floor computed offline from the teacher graph. Two build-time diagnostics accompany them: $D_{\mathrm{KL}}(p_{i,r}^T\|\mathcal{U}_{\mathrm{supp}})$ per scale, which detects a graph temperature that has flattened the targets into binary labels, and the pairwise cross-scale divergences $D_{\mathrm{KL}}(p_{i,r}^T\|p_{i,r'}^T)$, which detect scales that have become redundant. A near-zero $\mathrm{JS}_\omega$ is itself a warning: it certifies that the scales carry the same target and that, under a tied temperature, the multi-scale objective would be numerically identical to a single-scale one.

We also deduplicate the corpus before building the graph. Exact duplicate texts have $\cos=1$ by construction, so their logit is pinned at the ceiling with no gradient while they continue to absorb teacher probability mass from the candidates that could still move.

\begin{algorithm}[t]
\caption{HeatGeo Training}
\label{alg:heatgeo}
\begin{algorithmic}[1]
\REQUIRE Corpus $\mathcal{D}$, teacher $T$, student $S_{\theta}$, graph size $k$, scales $\bar{\mathcal{R}}=\{0\}\cup\mathcal{R}$, temperatures $\{\tau_r\}$, teacher temperature $\tau_t$, pool size $P$, candidate size $C$
\STATE Deduplicate $\mathcal{D}$; cache teacher embeddings $t_i=T(x_i)$
\STATE Build a mutual $k$NN teacher graph $W^T$
\STATE Compute $P^T=(D^T)^{-1}W^T$ and lazy walk $Q^T=\frac{1}{2}(I+P^T)$
\STATE Precompute diffusion pools $\mathcal{P}_i$ ($|\mathcal{P}_i|=P$) and targets $p_{i,r}^T$ from one walk snapshotted at each $r\in\mathcal{R}$; precompute same-source hard-negative pools $\mathcal{H}_i$
\STATE Compute the floor $\mathrm{JS}_{\omega}$ of \eqref{eq:floor} and the build-time target diagnostics
\FOR{each training epoch $e$}
    \STATE Redraw $C_i^{(e)}$ for every $i$: per-scale positives from $\mathcal{P}_i$, hard negatives from $\mathcal{H}_i$, random negatives
    \FOR{each mini-batch $B$}
        \STATE Encode anchors and candidate texts with $S_{\theta}$
        \STATE Form the shared column set $\bigcup_{i\in B}C_i^{(e)}$ and scatter the targets onto it
        \STATE Compute the direct target \eqref{eq:direct} on the shared columns from the cached $t_i$
        \STATE Compute $p_{i,r}^S$ at temperature $\tau_r$ for each $r\in\bar{\mathcal{R}}$
        \STATE Compute $\mathcal{L}_{\mathrm{HeatGeo}}$ and the excess over $\mathrm{JS}_{\omega}$
        \STATE Update $\theta$ by minimizing $\mathcal{L}_{\mathrm{HeatGeo}}$
    \ENDFOR
\ENDFOR
\RETURN Distilled student encoder $S_{\theta}$
\end{algorithmic}
\end{algorithm}
